\documentclass[11pt]{article}

\usepackage[margin=1.1in]{geometry}
\usepackage{amsmath,amssymb,amsthm,amscd}
\usepackage{booktabs}
\usepackage{microtype}
\usepackage{xcolor}
\usepackage[colorlinks=true,linkcolor=blue!55!black,citecolor=blue!55!black,urlcolor=blue!55!black]{hyperref}

% ---------- macros ----------
\newcommand{\cT}{\mathcal{T}}   % planar binary trees (free binary-tree structure)
\newcommand{\cW}{\mathcal{W}}   % words (free semigroup)
\newcommand{\cB}{\mathcal{B}}   % unordered binary trees (free commutative structure)
\newcommand{\cM}{\mathcal{M}}   % prime multisets (free commutative semigroup)
\newcommand{\cG}{\mathcal{G}}   % grouped quantities
\newcommand{\cA}{\mathcal{A}}   % flat quantities
\newcommand{\NN}{\mathbb{N}}
\newcommand{\pkg}[1]{\boxed{#1}}
\newcommand{\act}{\rho}
\newcommand{\physread}{\medskip\noindent\textbf{Physical reading.}\ }
\DeclareMathOperator{\supp}{supp}
\numberwithin{equation}{section}

\theoremstyle{plain}
\newtheorem{theorem}{Theorem}[section]
\newtheorem{proposition}[theorem]{Proposition}
\newtheorem{lemma}[theorem]{Lemma}
\newtheorem{corollary}[theorem]{Corollary}
\newtheorem{conjecture}[theorem]{Conjecture}
\theoremstyle{definition}
\newtheorem{definition}[theorem]{Definition}
\newtheorem{example}[theorem]{Example}
\newtheorem{question}[theorem]{Question}
\theoremstyle{remark}
\newtheorem{remark}[theorem]{Remark}

\title{Prime-Leaf Tree Theory\\[2mm]
\large Arithmetic as the statistical mechanics of multiplication histories}
\author{Giuseppe Barbalinardo}
\date{Working draft, \today}

\begin{document}
\maketitle

\begin{abstract}
We treat the integers as a gas with one mode per prime $p$, of energy $\log p$: the
setting of Julia's gas of numbers~\cite{julia} and of Bost--Connes. In the usual
reading, an integer
is a Fock state (its prime factorization is an occupation-number configuration) and
the Riemann zeta function is the free bosonic partition function. This paper studies
the same gas \emph{before} symmetrization and \emph{before} the interaction history
is forgotten. A microstate is then a \emph{multiplication tree}: an ordered,
parenthesized record of pairwise combinations, whose only conserved quantity is the
integer $N(T) = e^{E(T)}$ it evaluates to. Multiplication is not commutative or
associative at this level, because microstates remember their formation; ordinary
arithmetic is the fully coarse-grained description, and the Fundamental Theorem of
Arithmetic is precisely the statement that the occupation-number description is
faithful (zero residual entropy). Forgetting history and symmetrizing are independent
coarse-grainings, so there are four gases, each with an exact partition function
determined by the prime zeta function $P(\beta)$: the Riemann gas $\zeta(\beta)$, a
distinguishable-particle gas with inverse limiting temperature $1.39943\dots$, a
non-planar tree gas at $1.98847\dots$, and the full planar tree gas, whose
partition function obeys the self-consistency equation $Z = P + Z^2$ (a tree is a
prime or a pair of trees) and reaches its limiting temperature at
$\beta_H = \sigma^{*} = 2.59738\dots$, where $P(\sigma^{*}) = \tfrac14$. There the
partition function is finite and exactly $\tfrac12$, with the square-root branch point
universal for branched polymers: the transition is continuous, whereas the symmetrized
gases diverge at theirs. Associativity is, in this precise sense, a first-order
axiom. The complex-temperature zeros of the tree gas
are computed: they inherit the Riemann zeros through $P$, and their new branch points,
the solutions of $P(\beta) = \tfrac14$, do \emph{not} lie on a vertical line. Finally,
the bridge to addition: trees act as \emph{processes} on states of hierarchically
clustered matter. Under this action associativity becomes unobservable (processes
compose), commutativity fails at the level of states and survives only in expectation,
and distributivity is a theorem about the extensive observable rather than an axiom.
Primes themselves emerge as the cluster sizes that cannot crystallize into identical
cells. All combinatorial claims are verified by exhaustive enumeration and all
constants to high precision (Appendix~\ref{app:verification}).
\end{abstract}

\tableofcontents

%======================================================================
\section{Introduction}

\subsection{Multiplication before evaluation}

Ordinary arithmetic writes
\[
3 \times 2 = 6
\qquad\text{and}\qquad
1+1+1+1+1+1 = 6
\]
and treats the two results as the same kind of object. The starting intuition of this
paper is that they are not, and one liter of water is enough to see it. Pour six
one-liter bottles into a tank: that is addition, and it produces a flat six liters.
Now take $3 \times 2$ liters seriously as a construction: three cartons, each holding
two bottles. Take $2 \times 3$ liters: two cartons, each holding three. Both measure
six liters; neither \emph{is} six liters; and they are not each other. The slogan:

\begin{quote}
\emph{You never get six liters from multiplication. You get a structured object whose
shadow is $6$.}
\end{quote}

Physics has standard names for this situation: $6$ is a \emph{macrostate}, and the
packings that produce it are \emph{microstates}. We therefore refuse to let
multiplication collapse into numbers. A product is a \emph{multiplication tree}: the
recorded history of how the primes involved were combined, ordered and parenthesized,
unevaluated. (The cartons and bottles return in \S\ref{sec:bridge}, where the slogan
becomes a theorem.) The familiar integer is produced only by an evaluator, the shadow
map $N$, which plays the role of the measurement: it reads off the single conserved
quantity and discards the history. The classical laws of multiplication are then not properties of
multiplication at all. They are properties of the \emph{coarse-graining}: statements
about which distinct microstates the measurement fails to distinguish. In one line:

\begin{quote}
\emph{The primes are the particles; an integer is a measured energy; the Riemann
zeta function is a partition function; and ordinary arithmetic is what survives total
coarse-graining.}
\end{quote}

\subsection{The dictionary}\label{sec:dictionary}

The paper is written for readers who think in statistical mechanics. Every physical
term below is either an \emph{exact structural identity} (the object on the left
satisfies, verbatim, the defining equation of the object on the right) or is flagged
as imagery. Theorems remain theorems; the physics is load-bearing, not decorative.

\begin{table}[ht]
\centering\small
\begin{tabular}{@{}ll@{}}
\toprule
tree theory & physics \\
\midrule
prime $p$ & a mode of energy $\varepsilon_p = \log p$ \\
multiplication tree $T \in \cT$ & a composite: a recorded history of pairwise combinations \\
$\boxtimes$ (pair formation) & combination; on energies $E = \log N$ it is addition \\
shadow $N(T)$ & the measurement: the measured integer $e^{E(T)}$ \\
left--right order kept & ordered channels, as for planar diagrams \\
parenthesization kept & interaction history retained \\
quotient by associativity & forget the history, keep the particle list: asymptotic state \\
quotient by commutativity & Bose symmetrization \\
words $\cW$ & basis of the full (Boltzmann) Fock space \\
multisets $\cM$ & occupation numbers: the gas of numbers of Julia \\
degeneracy $g(n)$; \ $S(n) = \log g(n)$ & microstate count and Boltzmann entropy of the integer $n$ \\
Fund.\ Thm.\ of Arithmetic & the occupation-number description is faithful: $S_\cM \equiv 0$ \\
$Z(s) = \sum N^{-s}$, \ $s = \beta$ & canonical partition function at inverse temperature $\beta$ \\
$Z = P + Z^2$ & self-consistency equation (the Hagedorn--Frautschi bootstrap) \\
critical exponents $\sigma$ & inverse limiting (``Hagedorn'') temperatures $\beta_H$ \\
solutions of $P(s) = \tfrac14$ & complex-temperature zeros of $Z$ (its ``Fisher zeros'') \\
grouped quantities $\cG$ & states of hierarchically clustered matter \\
the action $\act(T)$ & a process applied to states; $\boxtimes \mapsto$ composition \\
flattening $\alpha$ & the extensive observable (particle number); expectation level \\
\bottomrule
\end{tabular}
\caption{The dictionary. The ``measurement'' and ``clustered matter'' rows are
imagery; the limiting temperature and complex-temperature zeros carry the standard
names (Hagedorn, Fisher) where the corresponding work is cited in the text. Every
other row is an exact identity.}
\label{tab:dictionary}
\end{table}

Everything is elementary and self-contained: finite combinatorics plus one-variable
complex analysis. No result depends on the imagery.

\subsection{Results}

\begin{enumerate}
\item \textbf{The microstate space} (\S\ref{sec:trees}). The world of trees is the
free binary-tree structure $\cT$ on the set of primes; the measurement is the unique
multiplication-preserving map $N\colon \cT \to (\NN_{\ge2},\cdot)$ with $N(L_p) = p$;
energies $E = \log N$ add under combination. The interacting world has no vacuum: $1$
is not the shadow of any tree, and the empty state exists only in the Fock
description.

\item \textbf{Coarse-graining is two-dimensional} (\S\ref{sec:square}). Forgetting the
interaction history (associativity) and symmetrizing (commutativity) are
\emph{independent} operations, giving a commuting square of four state spaces
$\cT, \cW, \cB, \cM$. The Fundamental Theorem of Arithmetic says exactly that the
final corner has one microstate per macrostate. Each level carries a computable
entropy: e.g.\ $S_\cT(12) = \log 6$. The degeneracies are Catalan $\times$ multinomial
for $\cT$, multinomial for $\cW$ (OEIS A008480), a sequence for $\cB$ that restricts
to the Wedderburn--Etherington numbers on prime powers and appears not to have been
recorded before, and $1$ for $\cM$. All verified exhaustively for $n \le 20000$.

\item \textbf{Four partition functions and a temperature ladder}
(\S\ref{sec:zeta}). All
four are exact functionals of the prime zeta function $P$: the free Bose gas
$\zeta(\beta)$ (the Euler product \emph{is} the textbook free-boson computation), a
geometric gas $P/(1-P)$, an Otter gas with argument-doubling equation, and the planar
tree gas obeying the self-consistency equation $Z = P + Z^2$. The inverse limiting
temperatures $\beta_H = 1/T_H$ form a strict ladder
\[
1 \;<\; 1.3994333287263303 \;<\; 1.9884768518009081 \;<\; \sigma^{*} = 2.5973851271346716 ,
\]
one rung per layer of remembered structure. At its limiting temperature the tree gas
converges with $Z_\cT(\sigma^{*}) = \tfrac12$ exactly and a square-root branch point:
the universal branched-polymer singularity, the same criticality that governs
planar-diagram counting in the Gaussian matrix model. The symmetrized gases instead
have poles: whether the limiting temperature is attainable is decided by
associativity. The classical Hagedorn--Frautschi--Nahm bootstrap, transplanted onto
the prime spectrum, slots into the same family with
$\beta_H = 2.1487511659\dots$ (Boltzmann counting) and $2.3072314551\dots$ (Bose
counting).

\item \textbf{Complex-temperature zeros} (\S\ref{sec:analytic}). $Z_\cT$ is algebraic
in $P$, so its complex-temperature singularities are inherited from $P$ (which
carries a branch point at $\beta = \rho/k$ for every nontrivial Riemann zero $\rho$,
and has the imaginary axis as a natural boundary), plus new branch points where
$P(\beta) = \tfrac14$. We compute the five with $0 < \operatorname{Im}\beta < 50$:
they recur near integer multiples of $2\pi/\log 2$ (the lightest mode, of energy
$\log 2$, sets the frequency) and do \emph{not}
lie on a vertical line. The naive tree analogue of the Riemann Hypothesis is false;
the right question is the distribution of these complex-temperature zeros, and that
question we answer: in every substrip of $\operatorname{Re} s > 1$ the tree zeros
have an exact asymptotic density (Theorem~\ref{thm:density}, via the
Jessen--Tornehave theory of almost periodic functions). Each zero occurs when
a single mode outshines the coherent rest, and the density is the total
dominance rate $\frac{1}{2\pi}\sum_p (\log p)\, q_p$. A census of all $59$ zeros to
height $500$ matches the formula's Monte Carlo prediction of $58.3$ to one percent.

\item \textbf{The additive bridge} (\S\ref{sec:bridge}). Trees act as
\emph{processes} on states of clustered matter (boxed heaps of unit quanta), and the
classical laws stratify by observability: associativity is invisible to the action
(exactly, because processes compose), while commutativity fails on states and
holds only in expectation, and distributivity holds only in expectation. ``You never
get six liters from multiplication'' becomes a theorem about the image of the action.
The primes themselves are then derivable inside the additive world: a size is
composite iff its heap can crystallize into $k \ge 2$ identical cells of $m \ge 2$
quanta; primes are the non-crystallizable sizes.

\item \textbf{Computational content} (\S\ref{sec:comp}). Two theorems mark what the
gas can and cannot compute. Microstate counting is \emph{shape-blind}: every
degeneracy depends only on the exponent profile of $n$, so no counting of histories
can locate primes: the no-free-lunch face of freeness. In the other direction, a
random integer seen at full resolution is a universal random object: a uniform
Catalan tree with a Gaussian, $\log\log x$, number of prime factors (Erd\H{o}s--Kac
lifted to trees). The interface where the theory genuinely meets the complexity of
finding primes is \emph{expression cost}: integer complexity is the first invariant
here that escapes $P$ (it distinguishes $4$ from $9$), Wilson's theorem makes
primality one factorial-expression away, and the open frontier, deterministic
prime generation, is a derandomization problem, mapped but not shortened.
\end{enumerate}

\subsection{Relation to existing work}\label{sec:priorart}

The Fock-level reading is classical: the gas of numbers of Julia~\cite{julia} and
Spector~\cite{spector}, deepened by Bost--Connes~\cite{bostconnes}, identifies
$\zeta(\beta)$ with a free bosonic partition function over the spectrum
$\{\log p\}$ and locates a limiting temperature at $\beta = 1$. The statistical
bootstrap and its limiting temperature are from~\cite{hagedorn,frautschi}, with the
analytical solution and the square-root criticality due to Nahm~\cite{nahm}.
Lee--Yang and Fisher zeros are from~\cite{leeyang,fisher}. On the mathematical side,
the free binary-tree structure on a set and its Catalan combinatorics are
standard~\cite{stanley}; the
degeneracy $g_\cT(n)$ appears as OEIS A144757 and the ordered-factorization count as
A008480~\cite{oeis}; ordered factorizations with arbitrary parts go back to
Kalm\'ar~\cite{kalmar}; the natural boundary of the prime zeta function is due to
Landau and Walfisz~\cite{landauwalfisz}. Closest in spirit to our refusal to
symmetrize is Loday's \emph{arithmetree}~\cite{loday,brunoyasaki}: arithmetic on
planar binary trees with non-commutative addition and one-sided distributivity.

Against this background the contributions here are: the two-dimensional
coarse-graining (the square, with the new non-planar corner, its Otter equation and
its limiting temperature); the exact critical value $Z_\cT(\sigma^{*}) = \tfrac12$
and the attainability dichotomy; the computation and non-collinearity of the
complex-temperature zeros; the process/state/expectation stratification of the
arithmetic laws; and the internal
derivation of the primes. One caution runs through the paper: a \emph{free} structure
carries no secrets about the primes themselves. Everything classically hard enters
through the single-particle spectrum, that is through $P(\beta)$, or through
the additive bridge. What the theory offers is a clean separation of the structural
from the arithmetic, and computable invariants in the separated layers.

%======================================================================
\section{Multiplication trees}\label{sec:trees}

\subsection{Modes and combination histories}

\begin{definition}[Multiplication trees]\label{def:trees}
Fix the set of prime leaves $\{L_p : p \text{ prime}\}$. The set $\cT$ of
\emph{multiplication trees} is defined by the grammar
\[
T \;::=\; L_p \;\mid\; (T_1 \boxtimes T_2).
\]
Thus $\cT$ is the \emph{free binary-tree structure} on the set of primes: $\boxtimes$
is ordered-pair formation, subject to no laws whatsoever.
\end{definition}

\physread $L_p$ is one particle in the mode of energy $\varepsilon_p = \log p$. A
tree is a composite that remembers how it was assembled: $T_1 \boxtimes T_2$ is the
combination of $T_1$ and $T_2$, in that order. Examples of \emph{distinct}
microstates:
\[
L_2, \qquad
L_2 \boxtimes L_3, \qquad
L_3 \boxtimes L_2, \qquad
(L_2 \boxtimes L_3) \boxtimes L_5, \qquad
L_2 \boxtimes (L_3 \boxtimes L_5).
\]
Non-commutativity and non-associativity are not axioms we impose; they are the
\emph{absence} of axioms. Keeping the left--right order is keeping the combination
channels ordered, as one does for planar diagrams of matrix-valued
fields~\cite{thooft}; keeping the parenthesization is keeping the interaction history:
which pair combined first. A multiplication tree is the parse tree of a
multiplication expression before evaluation; equivalently, a tree-level diagram built
from a single cubic vertex whose external legs are labeled by primes.

For $T \in \cT$ write $|T|$ for the number of leaves and $\mathrm{word}(T)$ for the
left-to-right sequence of leaf labels; $\supp(T)$ denotes the set of primes occurring.

\subsection{The measurement}

\begin{definition}[Shadow]\label{def:shadow}
The \emph{shadow map} $N \colon \cT \to \NN_{\ge 2}$ is defined recursively by
\[
N(L_p) = p, \qquad N(T_1 \boxtimes T_2) = N(T_1)\,N(T_2).
\]
Equivalently: $N$ is the unique multiplication-preserving map from $\cT$ to the
multiplicative semigroup $(\NN_{\ge2},\cdot)$ sending $L_p$ to $p$ (that is, the
unique map with $N(L_p)=p$ that respects $\boxtimes$).
\end{definition}

\physread Define the energy $E(T) = \log N(T)$. Then $E(L_p) = \varepsilon_p$ and
$E(T_1 \boxtimes T_2) = E(T_1) + E(T_2)$: energies add under combination
\emph{because} shadows multiply. The shadow is the measured value: the single
conserved
quantity, blind to history and to order. An integer is not a microstate; it is the
measured value $e^{E}$ shared by many microstates.

\begin{remark}[No vacuum]\label{rem:nounit}
$\cT$ has no unit object and $1$ is not the shadow of any tree: the image of
$N$ is exactly $\NN_{\ge 2}$. One could freely adjoin an empty tree $I$, but then
$I \boxtimes T \neq T$ (freeness again), which buys nothing. In gas language: the
interacting world has no vacuum state; the vacuum ($n = 1$, zero energy) exists only
in the Fock description, as the empty occupation configuration
(\S\ref{sec:square}).
\end{remark}

Composite numbers are not primitive objects of this world. There is no $L_6$; the
value $6$ occurs only as the common reading of the two trees $L_2\boxtimes L_3$
and $L_3\boxtimes L_2$. Likewise $8$ is the common reading of
$(L_2\boxtimes L_2)\boxtimes L_2$ and $L_2\boxtimes(L_2\boxtimes L_2)$: distinct
formation histories, one observable. Quantifying this degeneracy is the subject of the
next section.

%======================================================================
\section{Coarse-graining: from histories to Fock space}\label{sec:square}

\subsection{Two independent coarse-grainings}

\begin{definition}\label{def:congruences}
Let $\approx_a$ be the smallest congruence on $(\cT,\boxtimes)$ with
$(T_1\boxtimes T_2)\boxtimes T_3 \approx_a T_1\boxtimes(T_2\boxtimes T_3)$ for all
$T_i$, and $\approx_c$ the smallest congruence with
$T_1\boxtimes T_2 \approx_c T_2\boxtimes T_1$. Write
\[
\cW = \cT/\!\approx_a, \qquad
\cB = \cT/\!\approx_c, \qquad
\cM = \cT/(\approx_a \vee \approx_c).
\]
\end{definition}

Each quotient is a familiar free object: $\cW$ is the free semigroup on the primes
(nonempty \emph{words} in primes, i.e.\ ordered prime factorizations); $\cB$ is the
free \emph{commutative} binary-tree structure (trees with unordered branches); $\cM$
is the free commutative semigroup (nonempty finite \emph{multisets} of primes). The
shadow map
descends to all four corners, and the two identifications commute:
\[
\begin{CD}
\cT @>\text{forget history}>> \cW \\
@V\text{symmetrize}VV @VV\text{symmetrize}V \\
\cB @>>\text{forget history}> \cM
\end{CD}
\]
The shadow map of the final corner is where classical arithmetic begins:

\begin{proposition}[FTA as exactness of the collapse]\label{prop:fta}
The induced map $N \colon \cM \to \NN_{\ge2}$ is a bijection. This statement is
equivalent to the Fundamental Theorem of Arithmetic.
\end{proposition}

\begin{proof}
Elements of $\cM$ are nonempty prime multisets $\{p_1,\dots,p_k\}$ with
$N = p_1\cdots p_k$. Surjectivity of $N$ is the existence of a prime factorization;
injectivity is its uniqueness. \end{proof}

\physread The two arrows are the two things an observer can throw away.
\emph{Forgetting the history} keeps only the ordered list of constituent particles:
an asymptotic state, with the internal dynamics integrated out.
\emph{Symmetrizing} declares the particles indistinguishable: Bose
statistics. Doing both lands in the occupation-number description: an element of
$\cM$ is exactly a configuration $(\alpha_2, \alpha_3, \alpha_5, \dots)$, a basis
vector of the symmetric Fock space over the single-particle spectrum. $\cW$, the
half-way house that keeps order but forgets history, is a basis of the \emph{full}
(Boltzmann) Fock space. Proposition~\ref{prop:fta} then says: \emph{the
occupation-number description of the integers is faithful}. Unique factorization is
the absence of residual entropy at the Fock level.

\begin{remark}[The identifications are themselves structured]
The associativity identification is not a group action: two trees with the same word
are connected by a sequence of reassociations, and the set of parenthesizations of a
fixed word carries the \emph{Tamari lattice} structure, realized geometrically by the
associahedron~\cite{tamari}. Commutativity moves similarly generate the permutohedron.
The quotient square therefore has polytopal fibers, a refinement we do not pursue
here.
\end{remark}

\subsection{Degeneracies: the entropy of an integer}

For $n \ge 2$ with factorization $n = \prod_p p^{\alpha_p}$ write
$\Omega(n) = \sum_p \alpha_p$. For each corner $\mathcal{X}$ of the square let
\[
g_{\mathcal X}(n) = \#\{X \in \mathcal{X} : N(X) = n\},
\qquad
S_{\mathcal X}(n) = \log g_{\mathcal X}(n),
\]
the microstate count and the Boltzmann entropy of the macrostate $n$ at level
$\mathcal X$. Proposition~\ref{prop:fta} says $S_\cM \equiv 0$; the other three levels
have genuine entropy.

\begin{proposition}[Planar trees]\label{prop:gT}
\[
g_\cT(n) \;=\; C_{\Omega(n)-1}\;\frac{\Omega(n)!}{\prod_p \alpha_p!},
\qquad C_k = \frac{1}{k+1}\binom{2k}{k}.
\]
\end{proposition}

\begin{proof}
A tree with shadow $n$ has exactly $\Omega(n)$ leaves, by induction and
multiplicativity of $N$. The map $T \mapsto (\text{shape}(T), \mathrm{word}(T))$ is a
bijection from $\{T : N(T)=n\}$ to (binary tree shapes with $k=\Omega(n)$ leaves)
$\times$ (sequences of primes with product $n$). There are $C_{k-1}$
shapes~\cite{stanley} and $\Omega!/\prod \alpha_p!$ sequences (permutations of the
prime multiset). The two factors are independent: the fiber of the full collapse
$\cT \to \cM$ splits as histories $\times$ orderings.
\end{proof}

\begin{proposition}[Words]\label{prop:gW}
$g_\cW(n) = \Omega(n)!/\prod_p \alpha_p!$, with recursion
$g_\cW(n) = \sum_{p \mid n} g_\cW(n/p)$ and $g_\cW(p)=1$. This is OEIS A008480.
\end{proposition}

\begin{proposition}[Non-planar trees]\label{prop:gB}
$g_\cB$ satisfies $g_\cB(p) = 1$ and, for composite $n$,
\[
g_\cB(n) \;=\; \frac12 \sum_{\substack{d \mid n \\ 1<d<n}} g_\cB(d)\, g_\cB(n/d)
\;+\; \frac12\,[\,n \text{ is a square}\,]\; g_\cB(\sqrt n).
\]
On prime powers, $g_\cB(p^k)$ is the $k$-th Wedderburn--Etherington number
(OEIS A001190): 1, 1, 1, 2, 3, 6, 11, 23, \dots
\end{proposition}

\begin{proof}
An element of $\cB$ over composite $n$ is an unordered pair $\{B_1,B_2\}$ with
$N(B_1)N(B_2) = n$. Ordered pairs are counted by $\sum_{1<d<n} g_\cB(d)g_\cB(n/d)$;
the diagonal (possible only when $n$ is a square, both factors of shadow $\sqrt n$ and
equal as trees) is counted by $g_\cB(\sqrt n)$; unordered pairs with repetition are
(ordered $+$ diagonal)$/2$. On $p^k$ all leaves are identical and the count reduces to
unordered binary trees with $k$ indistinguishable leaves, the defining recursion of
the Wedderburn--Etherington numbers.
\end{proof}

\begin{remark}[A sequence not in the OEIS]\label{rem:newseq}
At the time of writing, the sequence $g_\cB(2), g_\cB(3), \dots =
1,1,1,1,1,1,1,1,1,1,2,\dots$ does not appear in the OEIS. The closest entry, A375120
(``complete binary unordered tree-factorizations''), uses a different convention: it
counts the two children as an \emph{ordered} pair whenever the factor split is
$(d,d)$, so it disagrees with Proposition~\ref{prop:gB} first at $n=144$ (it gives
$42$ where the true unordered count is $41$) and disagrees with
Wedderburn--Etherington on prime powers. Submitting $g_\cB$ is a small item of
independent interest.
\end{remark}

\begin{example}[The four levels over $n = 12$]\label{ex:twelve}
\[
\begin{array}{lll}
\cT: & 6 \text{ trees} &
(L_2\boxtimes L_2)\boxtimes L_3,\;
(L_2\boxtimes L_3)\boxtimes L_2,\;
(L_3\boxtimes L_2)\boxtimes L_2, \\
& & L_2\boxtimes(L_2\boxtimes L_3),\;
L_2\boxtimes(L_3\boxtimes L_2),\;
L_3\boxtimes(L_2\boxtimes L_2); \\[1mm]
\cW: & 3 \text{ words} & (2,2,3),\; (2,3,2),\; (3,2,2); \\[1mm]
\cB: & 2 \text{ unordered trees} & \{\{L_2,L_2\},L_3\},\; \{\{L_2,L_3\},L_2\}; \\[1mm]
\cM: & 1 \text{ multiset} & \{2,2,3\}.
\end{array}
\]
So the entropy of $12$ is $S_\cT(12) = \log 6$ at the history level,
$S_\cW(12) = \log 3$ for the particle list, $S_\cB(12) = \log 2$ non-planar, and $0$
in Fock space. And $g_\cT(12) = C_2 \cdot \tfrac{3!}{2!\,1!} = 2\cdot 3 = 6$, as predicted.
\end{example}

Table~\ref{tab:degeneracies} lists the three nontrivial counts for small $n$. All
formulas and recursions of this subsection were verified against brute-force
enumeration for every $n \le 20000$ (Appendix~\ref{app:verification}).

\begin{table}[ht]
\centering\small
\begin{tabular}{@{}rrrrr@{\hspace{9mm}}rrrrr@{}}
\toprule
$n$ & $\Omega$ & $g_\cT$ & $g_\cB$ & $g_\cW$ &
$n$ & $\Omega$ & $g_\cT$ & $g_\cB$ & $g_\cW$ \\
\midrule
 2 & 1 &  1 & 1 & 1 &  20 & 3 &  6 & 2 & 3 \\
 4 & 2 &  1 & 1 & 1 &  24 & 4 & 20 & 4 & 4 \\
 6 & 2 &  2 & 1 & 2 &  27 & 3 &  2 & 1 & 1 \\
 8 & 3 &  2 & 1 & 1 &  28 & 3 &  6 & 2 & 3 \\
10 & 2 &  2 & 1 & 2 &  30 & 3 & 12 & 3 & 6 \\
12 & 3 &  6 & 2 & 3 &  32 & 5 & 14 & 3 & 1 \\
16 & 4 &  5 & 2 & 1 &  36 & 4 & 30 & 6 & 6 \\
18 & 3 &  6 & 2 & 3 &  64 & 6 & 42 & 6 & 1 \\
\bottomrule
\end{tabular}
\caption{Degeneracies over $n$ at the three structured levels (primes and other rows
with all counts $1$ omitted). Note $g_\cT = C_{\Omega-1}\cdot g_\cW$ throughout, and
$g_\cB(2^k) = 1,1,1,2,3,6$ for $k=1,\dots,6$ (Wedderburn--Etherington).}
\label{tab:degeneracies}
\end{table}

%======================================================================
\section{Partition functions: self-consistency equations and the temperature ladder}
\label{sec:zeta}

\subsection{Four gases, four functional equations}

Write $E(X) = \log N(X)$ and let $\beta = s$ be the inverse temperature. Let
$P(s) = \sum_p p^{-s} = \sum_p e^{-s\varepsilon_p}$ denote the prime zeta function,
physically the single-particle partition function of the mode spectrum
$\{\log p\}$. For each corner
$\mathcal{X} \in \{\cT, \cW, \cB, \cM\}$ define the canonical partition function
\[
Z_{\mathcal X}(s) \;=\; \sum_{X \in \mathcal X} e^{-s E(X)}
\;=\; \sum_{X \in \mathcal X} \frac{1}{N(X)^{s}}
\;=\; \sum_{n \ge 2} \frac{g_{\mathcal X}(n)}{n^{s}}
\]
(for $\cM$ we include the empty configuration, of shadow $1$, so that
$Z_\cM = \zeta$ exactly).

\begin{theorem}[Functional equations]\label{thm:funceq}
As identities of Dirichlet series,
\begin{align}
Z_\cT(s) &= P(s) + Z_\cT(s)^2,
& Z_\cT &= \frac{1 - \sqrt{1 - 4P(s)}}{2} = \sum_{k\ge1} C_{k-1}\,P(s)^k,
\label{eq:cat}\\
Z_\cW(s) &= P(s) + P(s)\,Z_\cW(s),
& Z_\cW &= \frac{P(s)}{1-P(s)} = \sum_{k \ge 1} P(s)^k,
\label{eq:geo}\\
Z_\cB(s) &= P(s) + \tfrac12\bigl(Z_\cB(s)^2 + Z_\cB(2s)\bigr),
& Z_\cB &= 1 - \sqrt{1 - 2P(s) - Z_\cB(2s)},
\label{eq:otter}\\
Z_\cM(s) &= \zeta(s) = \prod_p \frac{1}{1-p^{-s}},
& \zeta &= \exp\Bigl(\sum_{k \ge 1} \frac{P(ks)}{k}\Bigr).
\label{eq:euler}
\end{align}
In \eqref{eq:cat} and \eqref{eq:otter} the branch of the square root is fixed by
$Z \to 0$ as $\operatorname{Re} s \to \infty$.
\end{theorem}

\begin{proof}
\eqref{eq:cat}: a tree is a leaf or an ordered pair of trees, and $N^{-s}$ is
multiplicative over $\boxtimes$, so the sum over pairs factors as $Z_\cT^2$. Solving
the quadratic and expanding gives the Catalan series.
\eqref{eq:geo}: a word is a prime or a prime followed by a word.
\eqref{eq:otter}: an element of $\cB$ is a leaf or an unordered pair
$\{B_1,B_2\}$, possibly equal; ordered pairs contribute $Z_\cB^2$, the diagonal
contributes $\sum_B N(B)^{-2s} = Z_\cB(2s)$, and unordered pairs with repetition are
half their sum, the same P\'olya computation as in Proposition~\ref{prop:gB}.
\eqref{eq:euler}: independence of the occupation numbers gives the Euler product, and
$\log \zeta(s) = \sum_p \sum_{k\ge1} p^{-ks}/k = \sum_k P(ks)/k$.
\end{proof}

\physread Each equation is a named piece of physics.
\emph{(\ref{eq:euler}) is the free Bose gas}: mode $p$ with occupancy $k$ costs energy
$k \log p$, so
\[
Z_\cM(\beta) \;=\; \prod_p \sum_{k \ge 0} e^{-\beta k \log p}
\;=\; \prod_p \frac{1}{1 - p^{-\beta}} \;=\; \zeta(\beta):
\]
the Euler product \emph{is} the textbook free-boson partition function: the
gas of numbers of Julia~\cite{julia}.
\emph{(\ref{eq:geo}) is the distinguishable-particle gas}: ordered $k$-particle states
with no $1/k!$, the counting that produces the Gibbs paradox.
\emph{(\ref{eq:cat}) is a self-consistency equation}: a composite is either an input
particle or a combination of two composites, with energy exactly conserved, so
$Z = P + Z^2$. This is the statistical bootstrap of
Hagedorn--Frautschi~\cite{hagedorn,frautschi}, here with input spectrum the primes
and the combination ordered and binary.
\emph{(\ref{eq:otter}) is the unordered version}: unordered combination, with the
$Z(2s)$ term the standard exchange correction for the two branches being
\emph{identical}, the same mechanism that separates Bose from Boltzmann counting.
Finally, Remark~\ref{rem:species} places all four in one family.

\begin{remark}[One gas per statistics]\label{rem:species}
The four right-hand sides are the four classical constructions of the symbolic
method~\cite{flajolet} (sequence, multiset, planar binary tree, unordered binary
tree), each applied to the single-particle element $P$. Choosing a corner of the
square is choosing a \emph{statistics}: what to remember of order and of history.
There is one partition function per statistics, all exact functionals of
$P(s), P(2s), \dots$, and the Riemann zeta function is the Bose member of the family.
Note where the vacuum lives: only the multiset construction has an empty structure,
which is why $n=1$ belongs to $\zeta$ and to no other corner
(Remark~\ref{rem:nounit}).
\end{remark}

\subsection{The temperature ladder}

\begin{theorem}[Inverse limiting temperatures]\label{thm:ladder}
Each series has a finite abscissa of convergence $\sigma_{\mathcal X}$, determined by
$P$:
\[
\begin{array}{llll}
\cM: & \sigma_\cM = 1 & \text{(pole of } \zeta), & Z \to \infty; \\[0.5mm]
\cW: & P(\sigma_\cW) = 1, & \sigma_\cW = 1.3994333287263303\dots, & Z \to \infty; \\[0.5mm]
\cB: & 2P(\sigma_\cB) + Z_\cB(2\sigma_\cB) = 1, & \sigma_\cB = 1.9884768518009081\dots, & Z \to 1; \\[0.5mm]
\cT: & P(\sigma^{*}) = \tfrac14, & \sigma^{*} = 2.5973851271346716\dots, & Z \to \tfrac12,
\end{array}
\]
and $1 < \sigma_\cW < \sigma_\cB < \sigma^{*}$: every remembered layer of structure
strictly raises the critical inverse temperature.
\end{theorem}

\begin{proof}[Proof sketch]
All coefficients are nonnegative, so each abscissa is located by the real singularity.
For $\cW$, the geometric series diverges exactly when $P \ge 1$. For $\cT$, the
Catalan series $\sum C_{k-1}P^k$ has radius of convergence $\tfrac14$ in $P$; since
$C_{k-1}4^{-k} \sim k^{-3/2}/(4\sqrt{\pi})$ is summable, the series still converges at
$P = \tfrac14$, with value $\sum_k C_{k-1}4^{-k} = \tfrac12$. For $\cB$, equation
\eqref{eq:otter} defines $Z_\cB$ by descending the argument-doubling cascade
($Z_\cB(2s)$ lies deep in the convergent region), and the singularity occurs where the
discriminant $1 - 2P(s) - Z_\cB(2s)$ vanishes, at which point $Z_\cB = 1$: this is
Otter's criticality condition for unordered trees~\cite{otter,flajolet}. The numerical
values are computed to the stated precision in Appendix~\ref{app:verification}; the
strict ordering follows from $g_\cM \le g_\cW, g_\cB \le g_\cT$ together with strict
inequality on a positive-density set of $n$.
\end{proof}

\physread The mechanism is Hagedorn's~\cite{hagedorn}. The number of microstates at
energy $E = \log n$ grows like $e^{\sigma_{\mathcal X} E}$ (up to powers), so the
Boltzmann sum $\sum e^{(\sigma_{\mathcal X} - \beta)E}$ diverges for
$\beta < \sigma_{\mathcal X}$: the gas has a maximal temperature
$T_H = 1/\sigma_{\mathcal X}$, the \emph{limiting temperature} (called the
\emph{Hagedorn temperature} in the hadron setting). Remembering order and history
multiplies the microstate count at each energy, raises the entropy per unit energy,
and lowers the limiting temperature:

\begin{center}
\begin{tabular}{@{}lllll@{}}
\toprule
gas & states & $\beta_H$ & $T_H = 1/\beta_H$ & transition \\
\midrule
Bose & occupation numbers & $1$ & $1$ & pole \\
Boltzmann & ordered factorizations & $1.399433\dots$ & $0.714575\dots$ & pole \\
non-planar tree & unordered histories & $1.988477\dots$ & $0.502897\dots$ & cusp, $Z_\cB \to 1$ \\
planar tree & ordered histories & $2.597385\dots$ & $0.385003\dots$ & cusp, $Z_\cT \to \tfrac12$ \\
\bottomrule
\end{tabular}
\end{center}

\begin{remark}[Pole versus cusp: is the limiting temperature attainable?]
\label{rem:cusp}
The symmetrized (associative) gases diverge at their critical point: $\zeta$ and
$Z_\cW$ have simple poles ($Z_\cW$ because $P'(\sigma_\cW) = -1.7311562\dots \ne 0$).
The history-keeping (non-associative) gases converge at theirs: expanding
\eqref{eq:cat} at $\sigma^{*}$,
\[
Z_\cT(s) \;=\; \tfrac12 - \sqrt{\,|P'(\sigma^{*})|\,}\,(s-\sigma^{*})^{1/2} + O(s-\sigma^{*}),
\qquad \sqrt{|P'(\sigma^{*})|} = 0.4796402876\dots,
\]
a square-root branch point with finite value, and similarly $Z_\cB \to 1$ at
$\sigma_\cB$. The canonical ensemble exists \emph{at} $T_H$, but the mean energy
\[
U(\beta) = -\frac{\partial}{\partial \beta} \log Z_\cT(\beta)
\;\sim\; \frac{\sqrt{|P'(\sigma^{*})|}}{2\,Z_\cT(\sigma^{*})}\,(\beta - \sigma^{*})^{-1/2}
\;\to\; \infty
\]
diverges as $\beta \to \sigma^{*+}$: a continuous transition with divergent
susceptibility, rather than a divergent partition function. Whether the limiting
temperature is attainable, the classic question of the bootstrap
literature~\cite{nahm}, is decided here by a single algebraic choice:
\emph{associativity is a first-order axiom}. The square-root branch point itself is
the universal singularity of tree and branched-polymer generating functions
(susceptibility exponent $\tfrac12$)~\cite{adj}, the same criticality that governs
planar-diagram counting in the Gaussian matrix model, where the Catalan numbers arise
as the planar moments and $\tfrac14$ is the edge of the Wigner
semicircle~\cite{bipz}.
\end{remark}

\begin{remark}[The classical bootstrap on the prime spectrum]\label{rem:bootstrap}
The historical statistical bootstrap model combines in one shot: a composite is an
input particle \emph{or an unordered cluster of $k \ge 2$ composites},
Boltzmann-weighted by $1/k!$. Its generating equation is
$G = \varphi + (e^{G} - 1 - G)$, whose square-root
singularity, first remarked by Nahm~\cite{nahm}, occurs when the input function
reaches the universal value
\[
\varphi_c \;=\; 2\log 2 - 1 \;=\; 0.3862943611\dots,
\qquad G_c = \log 2 .
\]
Evaluated on the empirical hadron spectrum, this criticality condition is what fixed
the original $T_H \approx 160$\,MeV. Transplanted onto the prime spectrum
($\varphi = P$), it gives the inverse Hagedorn temperatures
\[
\beta_H^{\mathrm{Boltz}} = 2.1487511659\dots,
\qquad
\beta_H^{\mathrm{Bose}} = 2.3072314551\dots,
\]
the Bose value arising from full multiset statistics
($e^{G} \mapsto \exp\sum_{k\ge1} G(ks)/k$, with criticality when
$P(s) + 2\sum_{k \ge 2} G(ks)/k$ reaches $2\log2 - 1$; see
Appendix~\ref{app:verification}). These classical gases
slot into the same family as the four corners (same input spectrum, same
square-root mechanism) but sit off the ladder of Theorem~\ref{thm:ladder} because
their symmetry-factor weights are fractional rather than integer counts. In bootstrap
language, the tree world $\cT$ is a bootstrap model whose decay chains are fully
resolved: an integer's prime factorization is its ultimate decay inventory, and the
primes are the stable particles.
\end{remark}

\begin{remark}[Growth of the counting functions]\label{rem:growth}
By the Wiener--Ikehara theorem applied to $Z_\cW$ (legitimate: the line
$\operatorname{Re} s = \sigma_\cW$ lies in the half-plane of analyticity of $P$, and
$|P(\sigma_\cW + it)| < 1$ for $t \ne 0$ by Lemma~\ref{lem:strict} below),
\[
\sum_{n \le x} g_\cW(n) \;\sim\; \frac{x^{\sigma_\cW}}{\sigma_\cW\,|P'(\sigma_\cW)|}
\;\approx\; 0.4128\, x^{1.39943\dots}.
\]
For the tree level, transfer heuristics for the square-root singularity suggest
$\sum_{n \le x} g_\cT(n) \asymp x^{\sigma^{*}} (\log x)^{-3/2}$; making this rigorous
is Problem~\ref{prob:tauberian}. Note also the extreme values: along $n = 2^k$ one has
$g_\cT(n) = C_{k-1}$, which grows like $n^{2}/(\log n)^{3/2}$ up to constants: the
microstate count over a single observable can approach the \emph{square} of the
observable. For comparison, Kalm\'ar's classical problem~\cite{kalmar} of ordered
factorizations with arbitrary parts $\ge 2$ has Dirichlet series $1/(2 - \zeta(s))$
and growth exponent $\rho = 1.7286472389981836\dots$ with $\zeta(\rho) = 2$; it sits
between our Boltzmann and non-planar rungs but belongs to a different family (its
parts need not be stable particles).
\end{remark}

%======================================================================
\section{Complex-temperature zeros of the arithmetic gas}\label{sec:analytic}

Lee and Yang taught that the thermodynamics of a system is encoded in the complex
singularities of its partition function, in complex fugacity for them and in complex
temperature for Fisher~\cite{leeyang,fisher}; the latter are the \emph{Fisher zeros}.
We now locate the complex-temperature
singularities of the tree gas. Since $Z_\cT = \tfrac12(1 - \sqrt{1-4P})$ is an
\emph{algebraic} function of $P$, they come in exactly two kinds: those inherited from
the single-particle input $P$, and the new branch points where the discriminant
vanishes.

\subsection{What the gas inherits: the Riemann zeros}

The prime zeta function continues to the right half-plane by M\"obius inversion of
\eqref{eq:euler}:
\begin{equation}\label{eq:Pcont}
P(s) \;=\; \sum_{k \ge 1} \frac{\mu(k)}{k}\, \log \zeta(ks),
\qquad \operatorname{Re} s > 0,
\end{equation}
valid off the singular set. Equation \eqref{eq:Pcont} shows that $P$ has logarithmic
branch points at $s = 1/k$ (from the pole of $\zeta$) \emph{and at $s = \rho/k$ for
every nontrivial zero $\rho$ of $\zeta$}; these points cluster on the imaginary axis,
which is a natural boundary (Landau--Walfisz~\cite{landauwalfisz}). Two consequences:

\begin{enumerate}
\item Lifting arithmetic to trees does not dissolve the Riemann Hypothesis. The
zeros of $\zeta$ reappear, relocated, inside $Z_\cT$ as branch points at $\rho/k$.
Where the Hilbert--P\'olya program seeks a Hamiltonian whose \emph{spectrum} is the
Riemann zeros~\cite{berrykeating}, the arithmetic gas offers the complementary
picture: the zeros enter as singularities of the single-particle input, while the gas
contributes its own critical singularities on top, and those are computable.
\item No one-variable partition function of this family can contain more analytic
information than $P$ itself. Genuinely new invariants must either weight the
\emph{shape} of histories (two-variable partition functions,
Problem~\ref{prob:twovar}) or couple to the additive world (\S\ref{sec:bridge}).
\end{enumerate}

\subsection{The new singularities: tree zeros}

The new singularities of $Z_\cT$ are the solutions of
\begin{equation}\label{eq:treezero}
P(s) = \tfrac14,
\end{equation}
square-root branch points of $Z_\cT$: the complex-temperature zeros of the tree gas.
We call them \emph{tree zeros} (they are the zeros of the discriminant $1 - 4P$). On
the real axis there is exactly one, $\sigma^{*}$, since $P$ is strictly decreasing
there. Off the axis:

\begin{lemma}\label{lem:strict}
For $\sigma > 0$ where $P(\sigma)$ converges and $t \neq 0$:
$|P(\sigma + it)| < P(\sigma)$.
\end{lemma}

\begin{proof}
$|\sum_p p^{-\sigma} p^{-it}| \le \sum_p p^{-\sigma}$ with equality iff all phases
$t \log p$ agree modulo $2\pi$. For $t \ne 0$ this fails already for $p = 2, 3$ since
$\log 3/\log 2$ is irrational.
\end{proof}

\begin{proposition}\label{prop:strip}
Every non-real tree zero satisfies $\operatorname{Re} s < \sigma^{*}$.
\end{proposition}

\begin{proof}
If $\sigma \ge \sigma^{*}$ and $t \neq 0$ then
$|P(\sigma+it)| < P(\sigma) \le P(\sigma^{*}) = \tfrac14$.
\end{proof}

The tree zeros with $0 < t < 50$ are listed in Table~\ref{tab:zeros} (method in
Appendix~\ref{app:verification}).

\begin{table}[ht]
\centering
\begin{tabular}{@{}cllc@{}}
\toprule
$k$ & $\operatorname{Re} s$ & $\operatorname{Im} s$ & $t \big/ \tfrac{2\pi}{\log 2}$ \\
\midrule
1 & $1.54427014392$ & $10.0122461654$ & $1.105$ \\
2 & $2.14220467429$ & $17.8224730193$ & $1.966$ \\
3 & $2.24736156780$ & $27.6431435747$ & $3.050$ \\
4 & $2.15502047591$ & $35.5407136013$ & $3.921$ \\
5 & $2.39048394209$ & $45.5394484588$ & $5.024$ \\
\bottomrule
\end{tabular}
\caption{All solutions of $P(s) = \tfrac14$ with $1.4 < \operatorname{Re} s < 2.7$,
$0 < \operatorname{Im} s < 50$ (each accompanied by its complex conjugate). A coarser
sweep of $1.05 < \operatorname{Re} s < 1.4$ found no further solutions in this height
range.}
\label{tab:zeros}
\end{table}

Two empirical observations. First, the ordinates hug integer multiples of
$2\pi/\log 2 = 9.0647\dots$: the partition function is almost periodic in
$\operatorname{Im}\beta$, and the recurrence frequency is set by the \emph{lightest
mode}, $\varepsilon_2 = \log 2$, whose Boltzmann factor $2^{-\beta}$ dominates $P$;
the interference of $3^{-\beta}$ and later modes then shifts each singularity
(violently so for $k=1$). Second, and decisively: \emph{the real parts wander}. The
naive tree analogue of the Riemann Hypothesis (all zeros on a vertical
line) is simply false. The correct lifted question is distributional, and it has a
precise answer.

\subsection{The density of tree zeros}\label{sec:density}

Within the half-plane of almost periodicity the tree zeros have an exact asymptotic
density, computable from a random-phase model of the gas.

\begin{theorem}[Density of the complex-temperature zeros]\label{thm:density}
For $\sigma > 1$ let $X_\sigma = \sum_p p^{-\sigma} e^{i\theta_p}$ with $(\theta_p)$
independent and uniform on $[0, 2\pi)$, and define the \emph{Jensen function} and the
\emph{dominance probabilities}
\[
\Phi(\sigma) = \mathbb{E}\,\log\bigl|X_\sigma - \tfrac14\bigr|,
\qquad
q_p(\sigma) = \mathbb{P}\Bigl(\,\bigl|X_\sigma - p^{-\sigma}e^{i\theta_p} -
\tfrac14\bigr| < p^{-\sigma}\Bigr).
\]
\begin{enumerate}
\item[(i)] For every $\sigma > 1$,
$\frac1T \int_0^T \log|P(\sigma + it) - \tfrac14|\,dt \to \Phi(\sigma)$ as
$T \to \infty$; $\Phi$ is convex, and $\Phi(\sigma) = \log\tfrac14$ for
$\sigma \ge \sigma^{*}$.
\item[(ii)] For $1 < \sigma_1 < \sigma_2$, the number $N_{[\sigma_1,\sigma_2]}(T)$ of
tree zeros with real part in $[\sigma_1, \sigma_2]$ and $0 < t \le T$ satisfies
\[
\frac{N_{[\sigma_1,\sigma_2]}(T)}{T} \;\longrightarrow\;
\frac{\Phi'(\sigma_2) - \Phi'(\sigma_1)}{2\pi}.
\]
\item[(iii)] $\Phi'(\sigma) = -\sum_p (\log p)\, q_p(\sigma)$. In particular the
density of tree zeros to the right of the line $\operatorname{Re} s = \sigma_1$ is
$\frac{1}{2\pi}\sum_p (\log p)\, q_p(\sigma_1)$, and $q_p(\sigma) = 0$ for
$\sigma \ge \sigma^{*}$, reproving Proposition~\ref{prop:strip}.
\end{enumerate}
\end{theorem}

\begin{proof}[Proof sketch]
In closed substrips of $\operatorname{Re} s > 1$ the series $P(s) - \tfrac14$
converges absolutely and uniformly, so it is a uniformly almost periodic analytic
function with frequency set $\{\log p\}$, linearly independent over $\mathbb{Q}$
(once more unique factorization). By Kronecker--Weyl the flow
$t \mapsto (t \log p \bmod 2\pi)_p$ equidistributes on the torus, identifying time
averages with expectations over the phases $(\theta_p)$: this is Bohr and Jessen's
argument for $\log\zeta$~\cite{bohrjessen}, and gives (i). The value $\log\tfrac14$
for $\sigma \ge \sigma^{*}$ follows by integrating out the phases one at a time: if
$\sum_p p^{-\sigma} \le \tfrac14$ then $|{-\tfrac14} + \text{rest}| \ge p^{-\sigma}$
pointwise for each peeled mode, and the circular average of $\log|W + a e^{i\theta}|$
equals $\log|W|$ whenever $|W| \ge a$. Statement (ii) is the fundamental theorem of
Jessen--Tornehave on mean motions and zeros of almost periodic
functions~\cite{jessentornehave}. For (iii), differentiate under the expectation and
average over $\theta_p$ first: for fixed $W$ with $|W| \neq a$,
\[
\frac{1}{2\pi}\int_0^{2\pi} \frac{a e^{i\theta}}{W + a e^{i\theta}}\, d\theta
\;=\;
\begin{cases} 1 & |W| < a, \\ 0 & |W| > a, \end{cases}
\]
by residues. The interchanges are justified: the characteristic function of
$X_\sigma$ is $\prod_p J_0(p^{-\sigma}|\xi|)$, with infinitely many Bessel factors
each decaying like $|\xi|^{-1/2}$, so $X_\sigma$ has a smooth bounded density. Hence
$\mathbb{E}\,|X_\sigma - \tfrac14|^{-1}$ is finite, uniformly on compact
$\sigma$-intervals; this dominates the $\sigma$-derivative, whose modulus is at most
$\sum_p (\log p)\, p^{-\sigma} \big/ |X_\sigma - \tfrac14|$, and gives the boundary
event $|W| = a$ probability zero.
\end{proof}

\physread Fix $\sigma$ and let $t$ run: $P(\sigma + it) - \tfrac14$ is a chain of
clock hands, one per prime, the $p$-hand of length $p^{-\sigma}$ turning at angular
speed $\log p$, the whole chain anchored at $-\tfrac14$. A tree zero at height $t$ is
the event that the tip of the chain passes through the origin. Whenever one hand is
longer than the anchored sum of all the others, that hand drags the tip around the
origin and forces one zero per revolution: this is the event measured by $q_p$. At
$\sigma = 1.30$ the hands have lengths $2^{-1.3} = 0.41$, $3^{-1.3} = 0.24$,
$5^{-1.3} = 0.12$, \dots, and the Monte Carlo values are $q_2 = 0.624$,
$q_3 = 0.148$, $q_5 = 0.042$, $q_7 = 0.017$ (Appendix~\ref{app:verification}).
Weighted by speed, the lightest mode accounts for only $59\%$ of the
singularities, which is exactly why the observed density \emph{exceeds} the naive
one-per-period guess $\log 2/2\pi$. Chains of rotating arms are the classical
mean-motion problem of celestial mechanics, Lagrange's secular perturbations; the
Jessen--Tornehave theory invoked above is its solution.

The numbers confirm the theorem. A census to height $500$
(Appendix~\ref{app:verification}) finds exactly $59$ zeros with
$\operatorname{Re} s > 1.3$; the dominance formula predicts
$\frac{500}{2\pi}\sum_p (\log p)\, q_p(1.30) = 58.3$, agreement to one percent
between two independent computations (Newton root-finding on
$\texttt{primezeta}$ against a random-phase Monte Carlo). One-per-period
would give $55.2$. The running counts are $N(100), \dots, N(500) = 11, 23, 35, 47,
59$; the mean spacing is $8.42$ against $2\pi/\log 2 = 9.06$; individual spacings
range from $1.31$ to $10.54$. The left edge, meanwhile, drifts: a coarser sweep of
$1.05 < \operatorname{Re} s < 1.30$ finds at least eight further zeros below the
strip, with real parts as low as $1.0130$ (at $t \approx 470.2$). The real parts
approach $1$ as the height grows, where almost periodicity fails and the continuation
\eqref{eq:Pcont}, with its Riemann-zero branch points, takes over: what happens
there is Problem~\ref{prob:zeros}.

%======================================================================
\section{The additive bridge: states, processes, and classical observables}
\label{sec:bridge}

The additive world of the theory is flat: $\cA = \{\,n\,\ell : n \in \NN_{\ge1}\}$,
heaps of a unit quantum $\ell$ (``one liter''), with
$+\colon \cA \times \cA \to \cA$ ordinary accumulation. The question left open so far
is how $\cT$ and $\cA$ interact. The answer proposed here: they do not interact by a
map $\cT \to \cA$. \emph{Multiplication acts on quantities} (trees are
\emph{processes}, not amounts), and the output of the action is neither a tree nor
a flat quantity but a third kind of object: a state of clustered matter.

\subsection{States: grouped quantities}

\begin{definition}[Grouped quantities]\label{def:grouped}
Define \emph{packages} $t$ and \emph{quantities} $f$ by the mutual grammar
\[
t \;::=\; \ell \;\mid\; \pkg{f}\,,
\qquad\qquad
f \;::=\; t_1 \oplus t_2 \oplus \cdots \oplus t_m \quad (m \ge 1),
\]
where $\pkg{f}$ is a box containing the quantity $f$, and $\oplus$ is concatenation of
finite sequences of packages. Let $\cG$ denote the set of quantities. For $k \ge 2$
define the \emph{uniform grouping}
\[
k \odot f \;=\; \underbrace{\pkg{f} \oplus \pkg{f} \oplus \cdots \oplus \pkg{f}}_{k}.
\]
The \emph{flattening} $\alpha \colon \cG \to \cA$ counts atoms:
$\alpha(\ell) = \ell$, $\alpha(\pkg f) = \alpha(f)$,
$\alpha(f \oplus g) = \alpha(f) + \alpha(g)$. We identify $\cA$ with
$(\NN_{\ge1}, +)$ and write $\alpha(f) \in \NN_{\ge 1}$.
\end{definition}

\physread A quantity is matter with a cluster hierarchy: quanta in bottles,
bottles in cartons; nucleons in nuclei, nuclei in atoms. $\oplus$ places systems side
by side; $\alpha$ is the extensive observable, total particle number, blind to
packaging. $\cG$ refines $\cA$: a state remembers how it is packed. ($\oplus$ is
associative by construction, being concatenation; one may further quotient $\cG$ by
reordering of packages, and nothing below changes, as the proofs never use the order.)

\begin{definition}[The action]\label{def:action}
Define $\act \colon \cT \to \mathrm{Maps}(\cG, \cG)$ by
\[
\act(L_p)(f) = p \odot f,
\qquad
\act(T_1 \boxtimes T_2) = \act(T_1) \circ \act(T_2).
\]
\end{definition}

\begin{example}\label{ex:liters}
$\act(L_3)(\ell) = \pkg{\ell}\,\pkg{\ell}\,\pkg{\ell}$: three bottles of one liter.
Then
\[
\act(L_2 \boxtimes L_3)(\ell)
= 2 \odot \bigl(\pkg{\ell}\,\pkg{\ell}\,\pkg{\ell}\bigr)
= \pkg{\,\pkg{\ell}\,\pkg{\ell}\,\pkg{\ell}\,}\;\pkg{\,\pkg{\ell}\,\pkg{\ell}\,\pkg{\ell}\,}\,,
\]
two cartons each holding three bottles; whereas $\act(L_3 \boxtimes L_2)(\ell)$ is
three cartons each holding two bottles. Neither is ``six liters''; both flatten to
$6\ell$. The slogan of \S1 is now a statement about the image of $\act$:
multiplication outputs clustered states, never flat ones.
\end{example}

\subsection{The stratification theorems}

The three classical laws now separate into three different layers of observability.

\begin{theorem}[Associativity is additively invisible, exactly]\label{thm:assoc}
The action $\act$ factors through the associativity quotient $\cT \to \cW$, and the
induced map $\cW \to \mathrm{Maps}(\cG,\cG)$ is injective. Consequently the congruence
$\{(T, T') : \act(T) = \act(T')\}$ is precisely $\approx_a$.
\end{theorem}

\begin{proof}
\emph{Factoring:} $\mathrm{Maps}(\cG,\cG)$ is a monoid under composition, composition
is associative, and $\act$ sends $\boxtimes$ to $\circ$; hence
$\act((T_1 \boxtimes T_2)\boxtimes T_3) = \act(T_1)\circ\act(T_2)\circ\act(T_3)
= \act(T_1 \boxtimes(T_2 \boxtimes T_3))$, and $\act(T)$ depends only on
$\mathrm{word}(T)$.

\emph{Injectivity:} let $w = p_1 p_2 \cdots p_k$ and evaluate on the single quantum:
$F = \act(w)(\ell) = p_1 \odot \act(p_2\cdots p_k)(\ell)$ consists of exactly $p_1$
identical top-level packages whose common content is $\act(p_2 \cdots p_k)(\ell)$.
So from $F$ one reads off $p_1$ (the number of top-level packages) and recurses on the
content; the recursion terminates at the boxless quantum $\ell$. Distinct words
therefore give distinct maps.
\end{proof}

\physread Why does associativity feel inevitable, when nothing in \S\ref{sec:trees}
imposes it? Because \emph{processes compose}: any semantics that interprets
multiplication as something one \emph{does} (anything of the form
$\boxtimes \mapsto \circ$ into an endomorphism monoid) must identify at least
$\approx_a$, since composition of maps is associative whether we like it or not. Our
semantics identifies \emph{exactly} $\approx_a$: nothing more is lost. Interaction
history is unobservable in any operational reading of multiplication; to observe it,
one would have to make states out of the processes themselves (currying), which is
precisely to leave the additive world. By contrast, nothing forces the process monoid
to be commutative, and indeed it is not.

\begin{theorem}[Commutativity survives until expectation]\label{thm:shadowaction}
For all $T \in \cT$ and $f \in \cG$:
\[
\alpha\bigl(\act(T)(f)\bigr) \;=\; N(T)\cdot \alpha(f).
\]
Hence the induced action on flat quantities $\cA$ is multiplication by the shadow, and
two trees act identically on $\cA$ iff $N(T) = N(T')$, i.e.\ (by
Proposition~\ref{prop:fta}) iff $T \approx_{ac} T'$. In particular
$\act(L_2 \boxtimes L_3) \ne \act(L_3 \boxtimes L_2)$ as maps of $\cG$, yet both have
the same expectation: every observer reading only particle number sees
multiplication by $6$.
\end{theorem}

\begin{proof}
Induction: $\alpha(p \odot f) = p\,\alpha(f)$ since boxes are $\alpha$-transparent,
and $\alpha(\act(T_1\boxtimes T_2)(f)) = N(T_1)\,\alpha(\act(T_2)(f)) =
N(T_1)N(T_2)\,\alpha(f)$. The characterization of equal flat actions is unique
factorization. Distinctness on $\cG$ is Theorem~\ref{thm:assoc} (different words).
\end{proof}

The two observation layers thus recover the two axes of the coarse-graining square:
\emph{clustered} observation quotients by exactly associativity; \emph{extensive}
observation quotients by exactly associativity and commutativity, and the
``exactly'' of the second statement is one more disguise of the Fundamental Theorem
of Arithmetic.

\begin{theorem}[Distributivity holds only in expectation]\label{thm:distrib}
For every $T \in \cT$ and $f, g \in \cG$:
\[
\act(T)(f \oplus g) \;\neq\; \act(T)(f) \oplus \act(T)(g)
\quad\text{in } \cG,
\qquad\text{yet}\qquad
\alpha\bigl(\act(T)(f \oplus g)\bigr) \;=\;
\alpha\bigl(\act(T)(f) \oplus \act(T)(g)\bigr).
\]
\end{theorem}

\begin{proof}
Let $p_1$ be the first letter of $\mathrm{word}(T)$. The left state consists of $p_1$
top-level packages; the right state consists of $2p_1$; since $p_1 \ge 2$ these
differ. Flattening both sides gives $N(T)(\alpha f + \alpha g)$ by
Theorem~\ref{thm:shadowaction}.
\end{proof}

``Process the merged system'' and ``process the parts, then merge'' are genuinely
different protocols with the same extensive outcome: the informal claim of \S1,
now a theorem. Note also the built-in one-sidedness: there is no
$(f \oplus g) \cdot T$ at all, because states are not processes. It is striking that
in Loday's arithmetree~\cite{loday}, where both arguments \emph{are} trees,
distributivity likewise survives on one side only.

\begin{center}
\begin{tabular}{@{}lccc@{}}
\toprule
law & in $\cT$ (microstates) & on states $\cG$ & in expectation ($\alpha$) \\
\midrule
associativity & fails (free) & \textbf{holds} & holds \\
commutativity & fails (free) & fails & \textbf{holds} \\
distributivity & not formulable & fails & \textbf{holds} \\
\bottomrule
\end{tabular}
\end{center}

\subsection{Where the primes come from}\label{sec:primesfrom}

The theory so far imports its particle spectrum from classical arithmetic. The bridge
removes this dependency: the primes can be \emph{found} inside $\cG$.

\begin{definition}\label{def:groupable}
Call $n \in \NN_{\ge 2}$ \emph{groupable} if some quantity consisting of $k \ge 2$
identical packages, each containing $m \ge 2$ quanta, flattens to $n$, that is, if
$n = km$ with $k, m \ge 2$. Call $n$ \emph{prime} if it is not groupable.
\end{definition}

\physread A size is composite iff a heap of that size can crystallize into $k$
identical unit cells; the primes are the non-crystallizable sizes, the stable
particles of the decay chains of Remark~\ref{rem:bootstrap}. This is, of course, the
classical definition of primality; that is the point. It is expressible entirely
inside the clustered additive world, using only $\odot$ and $\alpha$, with no imported
number labels. One can therefore \emph{define} the particle spectrum as the set of
non-groupable heap sizes and generate $\cT$ freely on it; the numerical label of a
prime leaf is recovered as the observable $\alpha$ of its heap. The foundational question
``are prime leaves labeled by known primes, or can the labels be derived?'' is thus
answered: derived. What cannot be derived from the free multiplicative structure alone
(\S\ref{sec:priorart}) becomes derivable the moment the additive world and the action
are in place, because primality is an observational property of clustering.

%======================================================================
\section{Arithmetic functions on trees}\label{sec:functions}

Classical arithmetic functions pull back along $N$. For the Euler totient,
\[
\varphi_\cT(T) \;=\; N(T) \prod_{p \in \supp(T)} \Bigl(1 - \frac1p\Bigr)
\;=\; \varphi\bigl(N(T)\bigr)
\]
by Euler's product formula; e.g.\
$\varphi_\cT(L_2 \boxtimes L_3) = 6\cdot\tfrac12\cdot\tfrac23 = 2$.
But this pullback factors through the double quotient: it cannot distinguish
$L_2\boxtimes L_3$ from $L_3 \boxtimes L_2$. A genuinely history-sensitive totient
should satisfy: (i) a recursion in $\boxtimes$; (ii) compatibility, i.e.\ it descends
to $\varphi \circ N$ after the double quotient; (iii) sensitivity, i.e.\ it separates
some pair of trees with equal shadow. No canonical candidate is known to us; we record
this as Problem~\ref{prob:totient}. The same design brief applies to M\"obius $\mu$,
divisor functions, and to two-variable partition functions
$Z_\cT(s, q) = \sum_T q^{\,\mathrm{stat}(T)} N(T)^{-s}$ for shape statistics
$\mathrm{stat}$ (depth, right-branch count, Tamari rank): the interesting ones are
exactly those \emph{not} expressible through $P$ alone (\S\ref{sec:analytic}).

%======================================================================
\section{Computational content}\label{sec:comp}

If the primes are the particles of this gas, can the gas help us \emph{find} them?
The question deserves a precise answer, because the temptation to answer yes is
strong, and wrong in an instructive way. Two theorems mark the boundary: microstate
counting is blind to primality (\S\ref{sec:nofreelunch}), while the microstate view
describes exactly what a random integer looks like (\S\ref{sec:randomtree}). We
then identify the one interface where the theory genuinely touches the complexity of
primes, the cost of \emph{expressions} (\S\ref{sec:cost}), and close with a map
of what ``finding a prime'' actually costs and where the open frontier lies
(\S\ref{sec:costlandscape}).

\subsection{No free lunch: degeneracy is shape-blind}\label{sec:nofreelunch}

\begin{proposition}[No free lunch]\label{prop:nofreelunch}
Let $n \ge 2$ have factorization $\prod_p p^{\alpha_p}$.
\begin{enumerate}
\item[(i)] $\displaystyle [\,n \text{ prime}\,] \;=\; g_\cT(n) \;-\;
\sum_{\substack{d \mid n \\ 1<d<n}} g_\cT(d)\, g_\cT(n/d)$.
\item[(ii)] For every corner $\mathcal X$ of the square, $g_{\mathcal X}(n)$ depends
only on the exponent multiset $\{\alpha_p\}$, not on which primes carry the exponents.
\item[(iii)] Given the factorization of $n$, the values $g_\cT(n)$ and $g_\cW(n)$ are
computable in time polynomial in $\log n$, and $g_\cB(n)$ in time $n^{o(1)}$.
\end{enumerate}
\end{proposition}

\begin{proof}
(i) is the defining recursion of $g_\cT$ rearranged. (ii): for $\cT$ and $\cW$ by the
closed formulas of Propositions~\ref{prop:gT}--\ref{prop:gW}; for $\cB$ by induction
along the divisor lattice, which is a product of chains of lengths $\alpha_p$ and
hence determined by the exponents alone (the square condition reads ``all $\alpha_p$
even''); for $\cM$ trivially. (iii): the closed formulas involve factorials of
$\Omega(n) = O(\log n)$, hence numbers with $O(\log n \log\log n)$ bits; the $\cB$
recursion runs over the divisor lattice, of size $d(n) = n^{o(1)}$.
\end{proof}

\physread The degeneracies cannot even distinguish $12 = 2^2\cdot3$ from
$18 = 2\cdot3^2$ from $20 = 2^2\cdot 5$ (all have $g_\cT = 6$, $g_\cB = 2$,
$g_\cW = 3$): they see the occupation \emph{profile}, never the identities of the
occupied modes. Identity (i) does determine primality, but its right-hand side
consumes the divisor structure of $n$, precisely what one would be trying to find,
and by (ii)--(iii) the entire structural layer is cheap dressing over the
factorization. This is the computational face of the caution of \S\ref{sec:priorart}:
freeness means the tree layer is polynomial-time syntax over its input spectrum.
Counting histories, at any corner of the square, can never be the hard step; whatever
algorithmic leverage exists must come from outside the counting.

\subsection{The random tree}\label{sec:randomtree}

What the microstate view \emph{does} compute is the statistics of a typical integer
seen at full resolution.

\begin{theorem}[Random tree]\label{thm:randomtree}
Draw $n$ uniformly from $\{2, \dots, \lfloor x\rfloor\}$, then draw $T$ uniformly
from the fiber $N^{-1}(n) \subset \cT$.
\begin{enumerate}
\item[(i)] Conditionally on $n$, the shape and the leaf word of $T$ are independent,
the shape uniform on the $C_{\Omega(n)-1}$ binary shapes and the word uniform on the
$\Omega(n)!/\prod_p \alpha_p!$ orderings of the prime multiset.
\item[(ii)] As $x \to \infty$,
$\bigl(\Omega(n) - \log\log x\bigr)/\sqrt{\log\log x}$ converges in distribution to a
standard Gaussian.
\item[(iii)] Consequently, conditionally on $\Omega(n) = k$ the tree is exactly a
uniform random binary tree with $k$ leaves; in particular its expected height, given
$k$, is asymptotic to $2\sqrt{\pi k}$.
\end{enumerate}
\end{theorem}

\begin{proof}
(i) The bijection of Proposition~\ref{prop:gT} identifies the fiber with
(shapes) $\times$ (words); the uniform measure on a product set has independent
uniform marginals. (ii) is the Erd\H{o}s--Kac theorem~\cite{erdoskac} for
$\omega(n)$, transferred to $\Omega$: since
$\mathbb{E}[\Omega - \omega] = \frac1x\sum_p \sum_{j \ge 2} \lfloor x/p^j \rfloor
\le \sum_p \frac{1}{p(p-1)} = O(1)$ and $\Omega \ge \omega$, the difference is
$o(\sqrt{\log\log x})$ in probability. (iii) The first claim restates (i); the height
asymptotic for uniform binary trees is Flajolet--Odlyzko~\cite{flajoletodlyzko}.
\end{proof}

\physread At full resolution a random integer is a small branched polymer: a Gaussian
number of constituents concentrated at $\log\log x$ (for $x = 10^{100}$, five or six
prime factors), a shape drawn from the \emph{universal} Catalan ensemble carrying no
arithmetic information at all (Proposition~\ref{prop:nofreelunch}(ii)
again), and height of order $\sqrt{\log\log x}$. Erd\H{o}s--Kac becomes the central
limit theorem for the constituent number of the arithmetic gas, and the literature on
random Catalan trees (profiles, path lengths, continuum limits) imports wholesale
into arithmetic through (iii). Sampling is equally well organized: one can generate a
uniform random integer below $x$ \emph{together with its Fock state} in polynomial
expected time, using primality tests as an oracle (Bach~\cite{bach}, simplified by
Kalai~\cite{kalai}): the occupation-number description is efficiently samplable,
even though deterministically \emph{aiming} at a prime is not
known to be.

\subsection{Expression cost: the invariant that escapes $P$}\label{sec:cost}

Everything so far in this paper is a functional of the input spectrum: the partition
functions of \S\ref{sec:zeta} through $P(s), P(2s), \dots$ (\S\ref{sec:analytic}),
the degeneracies through the exponent profile
(Proposition~\ref{prop:nofreelunch}). Here is the first invariant that is not.
Extend the grammar of \S\ref{sec:trees} with addition, as anticipated by the bridge:
define \emph{expressions}
\[
E \;::=\; 1 \;\mid\; (E_1 \oplus E_2) \;\mid\; (E_1 \boxtimes E_2),
\]
with evaluator $\mathrm{ev}(1) = 1$,
$\mathrm{ev}(E_1 \oplus E_2) = \mathrm{ev}(E_1) + \mathrm{ev}(E_2)$,
$\mathrm{ev}(E_1 \boxtimes E_2) = \mathrm{ev}(E_1)\,\mathrm{ev}(E_2)$, and cost the
number of $1$-leaves. The minimum cost
\[
\|n\| \;=\; \min\{\,\mathrm{cost}(E) : \mathrm{ev}(E) = n\,\}
\]
is the classical \emph{integer complexity} of Mahler--Popken~\cite{mahlerpopken}
(OEIS A005245~\cite{oeis}), with $3\log_3 n \le \|n\| \le 3\log_2 n$ for $n > 1$;
allowing subexpression reuse (straight-line programs) gives the Shub--Smale measure
$\tau(n) = O(\log n)$~\cite{shubsmale}.

Cost sees what no degeneracy can: $\|4\| = 4$ while $\|9\| = 6$, although $4 = 2^2$
and $9 = 3^2$ have the same exponent profile. It is, structurally, a bridge quantity:
it measures how expensively an integer can be synthesized once additions are allowed
inside the build. And it is exactly where the known connections between the structure
of multiplication and the complexity of primes live:

\begin{proposition}[Wilson's theorem as an expression statement]\label{prop:wilson}
\begin{enumerate}
\item[(i)] $n \ge 2$ is prime if and only if $(n-1)! \equiv -1 \pmod n$.
\item[(ii)] If $(x, N) \mapsto x! \bmod N$ were computable in time polynomial in
$\log N$, then factoring would be in polynomial time.
\end{enumerate}
\end{proposition}

\begin{proof}
(i) is Wilson's theorem. (ii): $\gcd(x! \bmod N,\, N) = \gcd(x!,\, N)$, which equals
$1$ exactly when $x$ is smaller than the least prime factor of $N$ and exceeds $1$
from that point on; binary search over $x \in [1, N]$ therefore finds the least prime
factor with $O(\log N)$ evaluations, and recursion factors $N$ completely.
\end{proof}

The same object carries the deeper conjectures. If $n!$ had straight-line programs of
size polylogarithmic in $n$, factoring would have polynomial-size circuits (folklore;
see Blum--Cucker--Shub--Smale~\cite{bcss}), and the Shub--Smale $\tau$-conjecture
(a polynomial bound on the integer roots of a polynomial in terms of its program
size) implies $\mathrm{P}_{\mathbb C} \neq \mathrm{NP}_{\mathbb C}$~\cite{shubsmale}.
In the vocabulary of this paper: \emph{the computational hardness of arithmetic is a
statement about the additive--multiplicative interface of \S\ref{sec:bridge}, about
the cost of synthesizing specific integers with additions, and Wilson's
theorem says primality is one factorial-expression away.}

The theory also suggests its own cost observable. Define the \emph{factored cost} and
the \emph{addition dividend}
\[
A(n) \;=\; \sum_{p^a \parallel n} a\,\|p\|,
\qquad
\delta(n) \;=\; A(n) - \|n\| \;\ge\; 0:
\]
$A(n)$ is the cost of the cheapest expression that respects the prime factorization
(synthesize each prime additively, then multiply), and $\delta(n)$ is the saving that
additions \emph{above} the factorization can buy. The long-standing open question
whether $\|2^k\| = 2k$ is exactly the statement $\delta \equiv 0$ on powers of two;
and $\delta$ is not identically zero: $\delta(55) \ge 1$, since
$A(55) = \|5\| + \|11\| = 13$ while $55 = 2 \cdot 27 + 1$ gives $\|55\| \le 12$.
Nothing about the distribution of $\delta$ appears to be known; see
Problem~\ref{prob:cost}.

\subsection{What finding a prime costs}\label{sec:costlandscape}

Finally, the target itself, stated honestly. \emph{Recognizing} primes is solved:
PRIMES $\in$ P~\cite{aks}. \emph{Hitting} one randomly is easy: by the prime number
theorem a random $n$-bit integer is prime with probability $\sim 1/(n \log 2)$, so
sampling plus testing generates an $n$-bit prime in expected time
$\widetilde{O}(n^3)$~\cite{crandallpomerance}. The genuine open problem, the
Polymath~4 problem~\cite{polymath}, is \emph{deterministic} generation:

\begin{center}\small
\begin{tabular}{@{}lll@{}}
\toprule
task & best known & status \\
\midrule
test an $n$-bit integer & deterministic $\mathrm{poly}(n)$~\cite{aks} & solved \\
generate, randomized & $\widetilde O(n^3)$: sample and test & routine \\
generate, deterministic & $\sim 2^{n/2}$: analytic $\pi(x)$ methods~\cite{polymath} & open, even under RH \\
generate, pseudodeterministic & $\mathrm{poly}(n)$, infinitely many $n$~\cite{chenluors} & frontier \\
quantum search & $\sim 2^{0.2625\,n}$: Grover~\cite{grover} on~\cite{bhp} & quadratic-type gain \\
\bottomrule
\end{tabular}
\end{center}

Even the Riemann Hypothesis only bounds prime gaps by $O(\sqrt{x}\log x)$, leaving a
$2^{n/2}$ search; a Cram\'er-type gap bound $O(\log^2 x)$~\cite{cramer} would make
``scan forward and test'' polynomial. The two attack surfaces are exactly the two
places this paper has located all remaining hardness. \emph{Gaps} are quantitative
statements about the additive--multiplicative interface: \S\ref{sec:bridge} names the
interface but is not yet quantitative. \emph{Derandomization}: the primes are dense
and efficiently recognizable, so deterministic generation is a problem of
\emph{aiming} (pseudorandomness rather than number-theoretic search), and the
natural formal home of ``dense yet hard to aim at'' is expression cost
(\S\ref{sec:cost}): a deterministic generator is, in the end, a short expression that
evaluates to a prime. The honest summary: the theory contributes a sharp map of where
the hardness cannot be (Proposition~\ref{prop:nofreelunch}), well-posed cost
questions at the interface where it can be, and no shortcut: by its own
no-free-lunch proposition, none was available from counting alone.

%======================================================================
\section{Open problems}\label{sec:open}

\begin{question}[Tree-zero distribution: beyond the almost periodic region]
\label{prob:zeros}
Theorem~\ref{thm:density} settles the density of tree zeros in substrips of
$\operatorname{Re} s > 1$. Remaining: (a) the left edge: the census finds real
parts as low as $1.013$ by height $500$; do they accumulate at $1$? Do zeros of the
continued $P - \tfrac14$ exist with $\operatorname{Re} s \le 1$, where almost
periodicity fails and the Riemann zeros enter through \eqref{eq:Pcont}? (b) the total
density $\lim_{\sigma_1 \downarrow 1} \frac{1}{2\pi}\sum_p (\log p)\,
q_p(\sigma_1)$: expected finite (the tail should behave like
$\sum_p (\log p)\, p^{-2\sigma}$), but its value and the speed of the leftward drift
are open. (c) pair correlations and the gap law of the ordinates: the census gives
spacings from $1.31$ to $10.54$ around mean $8.42$; is there level repulsion?
(d) a Lee--Yang-type question: is there a natural deformation of the gas for which
the singularities align on a line? (e) the same program for the level sets
$P(s) = c$, $0 < c < 1$: the tree zeros are $c = \tfrac14$, the Boltzmann-gas poles
$c = 1$.
\end{question}

\begin{question}[Tauberian rigor]\label{prob:tauberian}
Prove $\sum_{n \le x} g_\cT(n) \sim c\, x^{\sigma^{*}} (\log x)^{-3/2}$ (a
Delange-type theorem for Dirichlet series with a square-root singularity), and its
analogue for $g_\cB$.
\end{question}

\begin{question}[The cascade of the non-planar gas]\label{prob:cascade}
The continuation
\[
Z_\cB(s) = 1 - \sqrt{1 - 2P(s) - Z_\cB(2s)}
\]
propagates each singularity of $Z_\cB$ at $s_0$ to potential singularities at
$s_0/2, s_0/4, \dots$ and mixes them with those of $P$. Map this self-similar singularity structure;
determine whether $Z_\cB$ has a natural boundary strictly to the right of
$\operatorname{Re} s = 0$.
\end{question}

\begin{question}[Beyond the single-particle input]\label{prob:twovar}
Every one-variable partition function of \S\ref{sec:zeta} is a functional of
$P(s), P(2s), \dots$ Find a natural two-variable $Z_\cT(s,q)$ (weighting a shape
statistic of the history) that is provably \emph{not} such a functional, and study its
phase diagram in $(s,q)$.
\end{question}

\begin{question}[History-sensitive totient]\label{prob:totient}
Construct $\varphi^{\mathrm{tree}}$ satisfying (i)--(iii) of \S\ref{sec:functions},
ideally with an inclusion--exclusion (M\"obius) interpretation over subtrees.
\end{question}

\begin{question}[The full expression calculus]\label{prob:calculus}
\S\ref{sec:bridge} lets processes act on states, and \S\ref{sec:cost} equips the
two-sorted expression grammar with a cost. What is missing is its \emph{rewriting
theory}: treat distributivity as a directed rewrite rule (an evaluation step) rather
than an equation, study normal forms and confluence, and relate them to Loday's grove
addition~\cite{loday,brunoyasaki}.
\end{question}

\begin{question}[The addition dividend]\label{prob:cost}
Determine the behavior of $\delta(n) = A(n) - \|n\|$ (\S\ref{sec:cost}). Is
$\delta \equiv 0$ on powers of two (equivalently, $\|2^k\| = 2k$)? Is $\delta$
unbounded? What is its average order? Is $\|n\|$ computable in time polynomial in
$\log n$? And the structured refinement: quantify the \emph{cost of remembering},
the minimal cost of an expression whose multiplicative skeleton is constrained to a
fixed corner of the coarse-graining square.
\end{question}

\begin{question}[Axiomatics of the bridge]\label{prob:axioms}
Characterize the pair $(\cG, \alpha)$ abstractly: for which ``clustering worlds'' does
Theorem~\ref{thm:assoc} hold with \emph{exactly} $\approx_a$, and is $(\cG,\alpha)$
universal among them?
\end{question}

%======================================================================
\appendix
\section{Computational verification}\label{app:verification}

All computations are reproducible from the accompanying \texttt{primeleaf} Python
package (see the repository \texttt{README} for installation):
\texttt{python -m primeleaf verify-combinatorics} reproduces the combinatorial
checks below, \texttt{verify-zeta} the identities and constants, and
\texttt{census} and \texttt{dominance} the zero census and the density check.

\subsection{Combinatorial checks}

For every $n \le 20000$: the closed formula of Proposition~\ref{prop:gT} agrees with a
direct recursive enumeration $g_\cT(n) = [n \in \mathbb{P}] + \sum_{d \mid n,\, 1<d<n}
g_\cT(d)\, g_\cT(n/d)$; the multinomial of Proposition~\ref{prop:gW} agrees with the
first-letter recursion; and the unordered recursion of Proposition~\ref{prop:gB}
reproduces the Wedderburn--Etherington numbers
$1,1,1,2,3,6,11,23,46,98,207,451,983,2179$ at $n = 2^k$, $k \le 14$, as well as the
hand-enumerated values $g_\cB(12)=2$, $g_\cB(24)=4$, $g_\cB(30)=3$. The OEIS near-miss
A375120 (which uses ordered pairs on the diagonal) first diverges at $n = 144$: $42$
against the true $41$; a control query confirmed the OEIS search interface finds
A375120 from its own terms while the $g_\cB$ terms return no match (July 2026).

\subsection{Analytic checks and constants}

With $g_\cT(n)$, $g_\cB(n)$ enumerated for $n \le 20000$, partial sums of the
Dirichlet series were compared with the closed/functional forms of
Theorem~\ref{thm:funceq}:

\begin{center}\small
\begin{tabular}{@{}llll@{}}
\toprule
series & $s$ & partial sum vs.\ closed form & difference \\
\midrule
$Z_\cT$ & $3$ & $0.225547637865\ldots$ vs.\ $0.225705704214\ldots$ & $-1.6\cdot10^{-4}$ (tail) \\
$Z_\cT$ & $4$ & $0.084059063652\ldots$ vs.\ $0.084059066410\ldots$ & $-2.8\cdot10^{-9}$ \\
$Z_\cT$ & $6$ & agreement & $-3.1\cdot10^{-18}$ \\
$Z_\cB$ & $3$ & agreement & $-4.0\cdot10^{-7}$ \\
$Z_\cB$ & $4$ & agreement & $-1.1\cdot10^{-11}$ \\
\bottomrule
\end{tabular}
\end{center}

\noindent
the residual differences being exactly the truncation tails expected from the critical
exponents. Constants (via \texttt{mpmath.primezeta} at 25--30 significant digits):
\[
\begin{aligned}
\sigma^{*} &= 2.59738512713467162\dots, &\quad P'(\sigma^{*}) &= -0.230054805474\dots,\\
\sigma_\cB &= 1.98847685180090810\dots, &\quad Z_\cB(\sigma_\cB) &= 1\ (\text{verified to } 10^{-6}),\\
\sigma_\cW &= 1.39943332872633032\dots, &\quad P'(\sigma_\cW) &= -1.73115620186\dots,\\
\rho_{\text{Kalm\'ar}} &= 1.72864723899818362\dots, &\quad P(2) &= 0.45224742004106550\dots
\end{aligned}
\]
For the classical bootstrap of Remark~\ref{rem:bootstrap} on the prime spectrum:
Boltzmann counting, $P(\beta_H) = 2\log2 - 1$, gives
$\beta_H = 2.14875116590273\dots$; Bose counting, solving
$P(s) + 2\sum_{k\ge2} G(ks)/k = 2\log 2 - 1$ with $G$ computed by the
argument-doubling cascade, gives $\beta_H = 2.30723145510135\dots$ with
$G(\beta_H) = 0.6646652572\dots$

\subsection{Tree zeros}

The zeros of Table~\ref{tab:zeros} were located by a grid scan of
$|P_{\le 20000}(s) - \tfrac14|$ (truncated prime sum, step $0.02$) over
$[1.40, 2.70] \times [0.3, 50]$, followed by Newton polishing on the full
$\texttt{primezeta}$ at 20 significant digits; each reported zero satisfies
$|P(s) - \tfrac14| < 10^{-15}$. A coarser sweep of $[1.05, 1.40) \times [0.5, 50]$
(step $0.05 \times 0.25$) found no additional candidates; its only sub-threshold dip
flowed to zero \#1 of the table.

\subsection{Zero census and the density check}

The census of \S\ref{sec:density} (command \texttt{python -m primeleaf census})
extends the scan to $[1.30, 2.72] \times [0.3, 500]$ with primes up to $10^5$ (grid
$0.02 \times 0.05$, seed threshold $0.06$, Newton polish at 18 significant digits,
residuals below $10^{-14}$): $59$ zeros, real parts in $(1.3149,\, 2.5396)$, listed
in \texttt{data/zeros\_t500.txt}. The running counts are
$N(100), \dots, N(500) = 11, 23, 35, 47, 59$; the one-per-period line
$T\log 2/2\pi$ would give $11.0, 22.1, 33.1, 44.1, 55.2$; the mean spacing is
$8.417$, with extremes $1.31$ and $10.54$. A
coarse sweep of $[1.05, 1.30) \times [0.5, 500]$ (steps $0.05 \times 0.25$, not
exhaustive) finds eight further zeros with real parts down to $1.01298$ (at
$t \approx 470.24$). The dominance probabilities at $\sigma = 1.30$ were estimated
from $4 \cdot 10^5$ random-phase samples over the primes up to $10^4$:
\[
q_2 = 0.6237, \quad q_3 = 0.1484, \quad q_5 = 0.0416, \quad q_7 = 0.0167, \quad
q_{11} = 0.0051, \quad q_{13} = 0.0031,
\]
giving predicted density $0.11659$ per unit height, i.e.\ $58.3$ zeros to height
$500$, against $59$ observed.

%======================================================================
\begin{thebibliography}{99}

\bibitem{aks}
M.~Agrawal, N.~Kayal, and N.~Saxena,
\emph{PRIMES is in P},
Ann. of Math. \textbf{160} (2004), 781--793.

\bibitem{adj}
J.~Ambj\o rn, B.~Durhuus, and T.~Jonsson,
\emph{Quantum Geometry: A Statistical Field Theory Approach},
Cambridge University Press, 1997.

\bibitem{berrykeating}
M.~V.~Berry and J.~P.~Keating,
\emph{The Riemann zeros and eigenvalue asymptotics},
SIAM Review \textbf{41} (1999), 236--266.

\bibitem{bipz}
E.~Br\'ezin, C.~Itzykson, G.~Parisi, and J.-B.~Zuber,
\emph{Planar diagrams},
Comm. Math. Phys. \textbf{59} (1978), 35--51.

\bibitem{bostconnes}
J.-B.~Bost and A.~Connes,
\emph{Hecke algebras, type III factors and phase transitions with spontaneous symmetry
breaking in number theory},
Selecta Math. (N.S.) \textbf{1} (1995), 411--457.

\bibitem{brunoyasaki}
A.~Bruno and D.~Yasaki,
\emph{The arithmetic of trees},
arXiv:0809.4448 (2008).

\bibitem{bach}
E.~Bach,
\emph{How to generate factored random numbers},
SIAM J. Comput. \textbf{17} (1988), 179--193.

\bibitem{bhp}
R.~C.~Baker, G.~Harman, and J.~Pintz,
\emph{The difference between consecutive primes, II},
Proc. London Math. Soc. \textbf{83} (2001), 532--562.

\bibitem{bcss}
L.~Blum, F.~Cucker, M.~Shub, and S.~Smale,
\emph{Complexity and Real Computation},
Springer, 1998.

\bibitem{bohrjessen}
H.~Bohr and B.~Jessen,
\emph{\"Uber die Werteverteilung der Riemannschen Zetafunktion},
Acta Math. \textbf{54} (1930), 1--35; \textbf{58} (1932), 1--55.

\bibitem{chenluors}
L.~Chen, Z.~Lu, I.~C.~Oliveira, H.~Ren, and R.~Santhanam,
\emph{Polynomial-time pseudodeterministic construction of primes},
Proc. 64th IEEE Symposium on Foundations of Computer Science (FOCS), 2023.

\bibitem{cramer}
H.~Cram\'er,
\emph{On the order of magnitude of the difference between consecutive prime numbers},
Acta Arith. \textbf{2} (1936), 23--46.

\bibitem{crandallpomerance}
R.~Crandall and C.~Pomerance,
\emph{Prime Numbers: A Computational Perspective},
2nd ed., Springer, 2005.

\bibitem{erdoskac}
P.~Erd\H{o}s and M.~Kac,
\emph{The Gaussian law of errors in the theory of additive number theoretic
functions},
Amer. J. Math. \textbf{62} (1940), 738--742.

\bibitem{fisher}
M.~E.~Fisher,
\emph{The nature of critical points},
in: Lectures in Theoretical Physics, vol.~VII~C,
University of Colorado Press, Boulder, 1965, pp.~1--159.

\bibitem{flajoletodlyzko}
P.~Flajolet and A.~Odlyzko,
\emph{The average height of binary trees and other simple trees},
J. Comput. System Sci. \textbf{25} (1982), 171--213.

\bibitem{flajolet}
P.~Flajolet and R.~Sedgewick,
\emph{Analytic Combinatorics},
Cambridge University Press, 2009.

\bibitem{frautschi}
S.~Frautschi,
\emph{Statistical bootstrap model of hadrons},
Phys. Rev. D \textbf{3} (1971), 2821--2834.

\bibitem{grover}
L.~K.~Grover,
\emph{A fast quantum mechanical algorithm for database search},
Proc. 28th ACM Symposium on Theory of Computing (STOC), 1996, pp.~212--219.

\bibitem{hagedorn}
R.~Hagedorn,
\emph{Statistical thermodynamics of strong interactions at high energies},
Nuovo Cimento Suppl. \textbf{3} (1965), 147--186.

\bibitem{jessentornehave}
B.~Jessen and H.~Tornehave,
\emph{Mean motions and zeros of almost periodic functions},
Acta Math. \textbf{77} (1945), 137--279.

\bibitem{julia}
B.~L.~Julia,
\emph{Statistical theory of numbers},
in: Number Theory and Physics (Les Houches, 1989), Springer Proceedings in Physics
\textbf{47}, Springer, 1990, pp.~276--293.

\bibitem{kalai}
A.~Kalai,
\emph{Generating random factored numbers, easily},
J. Cryptology \textbf{16} (2003), 287--289.

\bibitem{kalmar}
L.~Kalm\'ar,
\emph{\"Uber die mittlere Anzahl der Produktdarstellungen der Zahlen},
Acta Litt. Sci. Szeged \textbf{5} (1931), 95--107.

\bibitem{landauwalfisz}
E.~Landau and A.~Walfisz,
\emph{\"Uber die Nichtfortsetzbarkeit einiger durch Dirichletsche Reihen definierter
Funktionen},
Rend. Circ. Mat. Palermo \textbf{44} (1920), 82--86.

\bibitem{leeyang}
C.~N.~Yang and T.~D.~Lee,
\emph{Statistical theory of equations of state and phase transitions. I. Theory of
condensation}, Phys. Rev. \textbf{87} (1952), 404--409;
T.~D.~Lee and C.~N.~Yang, \emph{II. Lattice gas and Ising model}, ibid., 410--419.

\bibitem{loday}
J.-L.~Loday,
\emph{Arithmetree},
J. Algebra \textbf{267} (2003), 507--539; arXiv:math/0112034.

\bibitem{mahlerpopken}
K.~Mahler and J.~Popken,
\emph{On a maximum problem in arithmetic},
Nieuw Arch. Wiskd. (3) \textbf{1} (1953), 1--15.

\bibitem{nahm}
W.~Nahm,
\emph{Analytical solution of the statistical bootstrap model},
Nucl. Phys. B \textbf{45} (1972), 525--553.

\bibitem{otter}
R.~Otter,
\emph{The number of trees},
Ann. of Math. \textbf{49} (1948), 583--599.

\bibitem{polymath}
D.~H.~J.~Polymath,
\emph{Deterministic methods to find primes},
Math. Comp. \textbf{81} (2012), 1233--1246.

\bibitem{shubsmale}
M.~Shub and S.~Smale,
\emph{On the intractability of Hilbert's Nullstellensatz and an algebraic version of
``NP $\neq$ P?''},
Duke Math. J. \textbf{81} (1995), 47--54.

\bibitem{oeis}
N.~J.~A.~Sloane (ed.),
\emph{The On-Line Encyclopedia of Integer Sequences},
\url{https://oeis.org}. Entries A000108, A001190, A005245, A008480, A144757, A375120.

\bibitem{spector}
D.~Spector,
\emph{Supersymmetry and the M\"obius inversion function},
Comm. Math. Phys. \textbf{127} (1990), 239--252.

\bibitem{stanley}
R.~P.~Stanley,
\emph{Catalan Numbers},
Cambridge University Press, 2015.

\bibitem{tamari}
F.~M\"uller-Hoissen, J.~M.~Pallo, and J.~Stasheff (eds.),
\emph{Associahedra, Tamari Lattices and Related Structures},
Progress in Mathematics \textbf{299}, Birkh\"auser, 2012.

\bibitem{thooft}
G.~'t~Hooft,
\emph{A planar diagram theory for strong interactions},
Nucl. Phys. B \textbf{72} (1974), 461--473.

\end{thebibliography}

\end{document}
